\documentclass[../../../../main.tex]{subfiles}
\begin{document}

\begin{flushright}
\footnotesize\textit{【自動文字起こし・要確認】}
\end{flushright}

[解] ${}_n \mathrm{C}_k$ ($k=1,2,\cdots, n-1$)の最大公約数を$d$とする。まず、$d$は${}_n \mathrm{C}_1 = pq$をわり切ることが必要で、$d=1,p,q,pq$である。ここで、$q < n$から、
$$
{}_n \mathrm{C}_q = \frac{n!}{q!(n-q)!} = \frac{pq \cdot (pq-1) \cdots \{ (p-1)q + 1 \}}{q (q-1) \cdots 2 \cdot 1}
$$
で、この分母、分子共に$q$で1回しかわり切れないから、${}_n \mathrm{C}_q$は$q$でわり切れない。同様に${}_n \mathrm{C}_p$も$p$でわり切れない。\\
以上から$d=1$が必要で、この時 ${}_n \mathrm{C}_k \in \mathbb{Z}$から十分である。

\end{document}
